\section[Material endpoint: retained parameter and arithmetic jets]{Material endpoint with the original parameter retained}
\label{sec:endpoint-jets}

The polynomial map and its provenance are those of
Section~\ref{sec:inverse-fibre-heat}~\cite{tao2026jacobian,speyer2026jacobian,poplett2026alias}.
We calculate the original Euclidean operator, not a diffusion chosen on
the target. The coordinate-gradient and divergence conventions agree
with Lychev--Koifman, Sections~2.3.3--2.3.4, especially (2.26)
\cite{lychevKoifman2019}; all uses here are proved by the chain rule below.
The preceding cancellation at $b=0$ is not a cancellation at every $b$.

\subsection{The full scalar diffusion on every retained real branch}

Keep $s=F_1/2$, $b=F_2$, $c=F_3$ and the original coordinates $(x,y,w)$.
Fix a real $b$ and introduce only the abbreviations
\[
 \beta=b/4,\qquad s_*=b^2/32=\beta^2/2,
 \qquad \eta=r-b/4.
\]
Their inverses are $b=4\beta$, $r=\eta+b/4$; the original parameter
has not been set to zero. On $c=0$, the inverse chart of
\eqref{eq:mag-shifted-branch} is
\begin{equation}
 H_b(\eta)=\left(-\frac1{2\eta},\beta+3\eta,
                          26\eta^2+6\beta\eta\right),\qquad
 s(H_b(\eta))=s_*-\eta^2/2,\quad \eta\ne0.
 \label{eq:ej-chart}
\end{equation}
It satisfies $F(H_b(\eta))=(2s,b,0)$. The separate finite point
$(0,b,2s-4b^2)$ remains on that same fibre, including at $s=s_*$.
On either real branch, the original material time satisfies
\[
 \frac{d\eta}{d\tau}=-\eta^{-1},\qquad
 \eta(\tau)=\operatorname{sgn}(\eta_0)
                  \sqrt{\eta_0^2-2\tau},\qquad \tau<\eta_0^2/2.
\]
Indeed $Us=1$, $Ub=Uc=0$ on this slice, and differentiation of
$s=s_*-\eta^2/2$ proves the equation. Thus $H_b$ is an actual
escaping material orbit, not merely a curve on a fixed fibre.

Here are all three gradient entries and the full Laplacian, evaluated
only after differentiation in the original three spatial variables:
\begin{align}
 g_x(\eta)&=\frac{(\beta+3\eta)
 (3\beta^3+17\beta^2\eta+9\beta\eta^2+3\eta^3)}{4\eta},\notag\\
 g_y(\eta)&=\frac{3\beta^3+9\beta^2\eta+17\beta\eta^2+3\eta^3}
                         {8\eta^2},\notag\\
 g_w(\eta)&=-\frac{(\beta+\eta)^3}{16\eta^3},\notag\\
 L_b(\eta)&=\frac{9\beta^2}{4\eta^2}+\frac{9\beta}{2\eta}
 +\frac{13}{4}-6\beta^4-66\beta^3\eta-246\beta^2\eta^2
                 -342\beta\eta^3-108\eta^4,\notag\\
 N_b(\eta)&=g_x(\eta)^2+g_y(\eta)^2+g_w(\eta)^2.
 \label{eq:ej-full-coefficients}
\end{align}
Here $g_j=(\partial_j s)\circ H_b$, $L_b=(\Delta_Es)\circ H_b$,
and $N_b=|\nabla_Es|^2\circ H_b$ for real variables. These three
displayed squares specify every coefficient of $N_b$, including
all regular terms. For holomorphic continuation the same algebraic
sum of squares is used; it is not called a Hermitian norm there.

To verify the formulas, set $v=1+xy$ in
$s=(v^3w+y^2(v+3v^2))/2$. Differentiate using $v_x=y$, $v_y=x$,
$v_w=0$ as in \eqref{eq:mag-Euclidean-diffusion}. Substitute
$v(H_b)=-(\beta+\eta)/(2\eta)$, together with
\eqref{eq:ej-chart}, into those three first derivatives and the
displayed polynomial formula for $\Delta_Es$ in that section.
Collecting powers gives precisely \eqref{eq:ej-full-coefficients}.
The scalar chain rule, before restriction, now proves for every
$C^2$ function $f$ near $s_*$
\begin{equation}
 (\Delta_E(f\circ s))(H_b(\eta))
 =N_b(\eta)f''(s_*-\eta^2/2)+L_b(\eta)f'(s_*-\eta^2/2).
 \label{eq:ej-diffusion}
\end{equation}
This is a map from the scalar function and its two jets to a function
on the retained chart; no fibre-independent scalar generator is inferred.

For $b\ne0$, the exact leading terms are
\begin{equation}
 \eta^3(g_x,g_y,g_w)\longrightarrow(0,0,-\beta^3/16),\qquad
 \eta^6N_b\longrightarrow\beta^6/256,\qquad
 \eta^2L_b\longrightarrow9\beta^2/4.
 \label{eq:ej-leading}
\end{equation}
They follow by multiplying the rational expressions by the stated
powers and evaluating at zero. For instance, along either branch
$\nu N_b\sim\nu\beta^6/(2048(T-\tau)^3)$, with the original
viscosity $\nu>0$ and $T=\eta_0^2/2$. For $b=0$, instead,
\[
 (g_x,g_y,g_w)=(9\eta^3/4,3\eta/8,-1/16),\qquad
 N_0=81\eta^6/16+9\eta^2/64+1/256,\quad L_0=13/4-108\eta^4.
\]
Taking $\eta=-z/2$ recovers every coefficient of
\eqref{eq:mag-orbit-diffusion}. This explicitly locates the earlier
finite limit at the special value $b=0$.

\subsection{Exact analytic-jet classification of bounded diffusion}

\begin{theorem}\label{thm:ej-jets}
Let $b\in\mathbb R\setminus\{0\}$, and let $f$ be complex analytic
near $s_*=b^2/32$, with its actual convergent expansion
$f(s_*+h)=\sum_{j\ge0}a_jh^j$.
Along either real branch of \eqref{eq:ej-chart}:
\begin{enumerate}
 \item The physical gradient $\nabla_E(f\circ s)$ is bounded exactly
 when $a_1=a_2=0$. In that case it tends to zero.
 \item The full Euclidean Laplacian in \eqref{eq:ej-diffusion} is
 bounded exactly when $a_1=a_2=a_3=a_4=0$. Its limit is then
 $-5\beta^6a_5/512$.
\end{enumerate}
No condition is imposed on $a_0$ in either statement.
\end{theorem}
\begin{proof}
The gradient is exactly $f'(s_*-\eta^2/2)(g_x,g_y,g_w)$.
Its last component has successive terms
$-\beta^3a_1/(16\eta^3)$ and, once $a_1=0$,
$\beta^3a_2/(16\eta)$. Their coefficients are nonzero when
the indicated $a_j$ is nonzero, so boundedness forces both to vanish.
Conversely $f'=O(\eta^4)$ if they vanish; the gradient entries
are $O(\eta^{-3})$, proving convergence to zero.

For the second statement it suffices to retain the first six
potentially negative Laurent powers. A convergent Taylor remainder
after $a_5h^5$ has second derivative $O(h^4)=O(\eta^8)$ and first
derivative $O(h^5)$; multiplying by $N_b=O(\eta^{-6})$ and
$L_b=O(\eta^{-2})$ leaves no pole and gives a vanishing remainder.
The coefficients of $\eta^{-6}$ and $\eta^{-5}$ are
\[
 \frac{\beta^6a_2}{128},\qquad \frac{3\beta^5a_2}{64}.
\]
Boundedness therefore forces $a_2=0$. After that substitution,
the coefficients of $\eta^{-4}$ and $\eta^{-3}$ are
\[
 -\frac{3\beta^6a_3}{256},\qquad
 -\frac{9\beta^5a_3}{128},
\]
so $a_3=0$. After both substitutions, the remaining pole coefficients
are exactly
\begin{equation}
 [\eta^{-2}]\Delta_E(f\circ s)
   =\frac{3\beta^2}{256}(192a_1+\beta^4a_4),\qquad
 [\eta^{-1}]\Delta_E(f\circ s)
   =\frac{9\beta}{128}(64a_1+\beta^4a_4).
 \label{eq:ej-final-poles}
\end{equation}
All these coefficients follow by expanding the three squares in
\eqref{eq:ej-full-coefficients} and
$f'=\sum ja_j(-\eta^2/2)^{j-1}$,
$f''=\sum j(j-1)a_j(-\eta^2/2)^{j-2}$.
Subtracting the two zero equations in parentheses gives $128a_1=0$;
then $a_4=0$, since $\beta\ne0$. A finite Laurent principal part
cannot be bounded on a real punctured half-interval unless every
negative-power coefficient vanishes: multiply by its highest pole
power and take the limit. This proves necessity on either branch.

If $a_1$ through $a_4$ vanish, then
$f''=-(5/2)a_5\eta^6+O(\eta^8)$ and $f'=O(\eta^8)$.
Equation~\eqref{eq:ej-leading} proves both sufficiency and the limit
$-(5/2)a_5\beta^6/256$. This also checks all possible cancellations
between the two terms of the actual Euclidean operator.
\end{proof}

In particular this theorem applies to the actual $Z_t(s)=8H_t(2s)$
at each fixed real heat time, with
$a_j=Z_t^{(j)}(s_*)/j!$. It does not replace $Z_t$ by its Taylor
polynomial: the proof explicitly controls the analytic remainder.
At an existing zero of order $m\ge2$, the first nonzero diffusion
term is
\[
 \frac{\beta^6}{256}m(m-1)a_m(-1/2)^{m-2}\eta^{2m-10}.
\]
For $m\ge5$ it is bounded; for $m=2,3,4$ it diverges. At a simple
zero it also diverges, by the complete pole classification, regardless
of possible cancellations involving $a_4$. These statements concern
the Euclidean derivative of the pulled-back function. A zero is not
created by this calculation, and its location has not been presumed.

\subsection{All fixed covectors at the original $b=0$ trajectory}

The scalar cancellation and the full material tensor can now be
compared without discarding transverse covectors. Put
$q_0=(1,-3/2,13/2)$, $z=\sqrt{1-8\tau}>0$, and
$\gamma=(z^{-1},-3z/2,13z^2/2)$ as before. Direct differentiation of
the original $F$ gives
\[
 J_0=\begin{pmatrix}-9/16&-3/8&-1/8\\3/8&25/4&3/4\\-17/2&-3&-1\end{pmatrix},
 \quad
 J_\gamma=\begin{pmatrix}-9z^3/16&-3z/8&-1/8\\3z^2/8&25/4&3/(4z)\\-17/2&-3/z^2&-1/z^3\end{pmatrix}.
\]
The full material derivative and inverse are
$A=J_\gamma^{-1}BJ_0$, $C=J_0^{-1}B^{-1}J_\gamma$,
where $B=I+3(1-z^2)E_{23}/4$. Differentiating
$F\circ X_\tau=(F_1+2\tau,F_2+6\tau F_3,F_3)$ proves these
identities for all initial directions, not just the orbit tangent.
Both determinants of $J$ are $-2$ and $\det B=1$, so $C=A^{-1}$
and $\det C=1$.

For an arbitrary fixed real initial covector
$\ell=(\ell_1,\ell_2,\ell_3)^{\mathsf T}$, define
\[
 \lambda=\frac{2\ell_1+3\ell_2-18\ell_3}{8},\qquad
 \mu=\ell_2-3\ell_3,\qquad \kappa=\ell_3-\lambda.
\]
These are coordinates on the entire covector space, with inverse
\[
 \ell_1=(17\lambda-3\mu+9\kappa)/2,\quad
 \ell_2=3\lambda+\mu+3\kappa,\quad \ell_3=\lambda+\kappa.
\]
Matrix multiplication in the displayed factorizations gives
\begin{equation}
 C^{\mathsf T}\ell=
 \left(\frac{17\lambda}{2}-\frac{3\mu z^2}{2}
                         +\frac{9\kappa z^3}{2},\quad
       \frac{3\lambda}{z^2}+\mu+3\kappa z,\quad
       \frac{\lambda}{z^3}+\kappa\right)^{\mathsf T}.
 \label{eq:ej-all-covectors}
\end{equation}
One can check every entry without inverting a symbolic flow: substitute
the inverse coordinates for $\ell$ into $J_0^{-\mathsf T}\ell$,
apply $B^{-\mathsf T}$, and then apply $J_\gamma^{\mathsf T}$.
This proves the formula for arbitrary $\ell$.

It follows that $C^{\mathsf T}\ell$ is bounded exactly on the plane
$2\ell_1+3\ell_2-18\ell_3=0$. For $\lambda\ne0$,
$z^3C^{\mathsf T}\ell\to\lambda e_3$ and
$z^6\ell^{\mathsf T}CC^{\mathsf T}\ell\to\lambda^2$.
For $\lambda=0$, its complete diffusion cost is
\begin{equation}
 \nu\ell^{\mathsf T}CC^{\mathsf T}\ell
 =\nu\left\{\frac94z^4(-\mu+3\kappa z)^2
                    +(\mu+3\kappa z)^2+\kappa^2\right\}
 \longrightarrow\nu(\mu^2+\kappa^2).
 \label{eq:ej-plane-energy}
\end{equation}
The limit is positive for every nonzero covector on that plane,
by the explicit inverse coordinate map. There is no additional
intermediate pole order for fixed covectors: setting $\lambda=0$
cancels both negative powers in \eqref{eq:ej-all-covectors}.
Time-dependent covectors are not covered by this fixed-input assertion.

The actual initial arithmetic covector is
$ds(q_0)=(-9/32,-3/16,-1/16)$, giving
$(\lambda,\mu,\kappa)=(0,0,-1/16)$. Formula~\eqref{eq:ej-all-covectors}
recovers exactly $ds(\gamma)=(-9z^3/32,-3z/16,-1/16)$ and
the finite cost $\nu/256$. The generic example $\ell=e_3$ has
$\lambda=-9/4$ and cost asymptotic $81\nu z^{-6}/16$.
Thus the two behaviors are obtained from one full inverse transpose.

\subsection{Exact compact Weil-test channels at a shifted branch}

Retain $\mathcal E=C_c^\infty(\mathbb R;\mathbb C)$,
$\widehat g(s)=\int g(u)e^{-isu}\,du$ and $T=-i\partial_u$.
Set $\sigma=s_*$ for this subsection. Define the actual endpoint
test subspace
\[
 \mathcal B_\sigma=\{g\in\mathcal E:
                  \widehat g^{(j)}(\sigma)=0,\ 1\le j\le4\}.
\]
By Theorem~\ref{thm:ej-jets} this is exactly the compact test space
whose Fourier transform has bounded pulled-back Euclidean Laplacian
at this nonzero-$b$ branch. Notice that its zeroth Fourier value
is unrestricted. Equivalently its four moments
$\int u^j e^{-i\sigma u}g(u)\,du$, $1\le j\le4$, vanish.

Here is a complete constructive parametrization, including that
unrestricted value. Choose any real smooth bump $\chi$ supported
in an interval $[a,d]$, $a<d$, with $\int\chi\ne0$, and put
$\chi_\sigma(u)=e^{i\sigma u}\chi(u)$.
For example, on $(a,d)$ one may take
$\chi(u)=\exp(-1/((u-a)(d-u)))$, extended by zero; every derivative
vanishes at each endpoint since an exponential dominates every
inverse power, so this is a positive smooth compact bump.
Then $\widehat\chi_\sigma(\sigma)=\int\chi\ne0$.
Let $P_4(s)$ be the degree-four Taylor polynomial at $\sigma$ of
$1/\widehat\chi_\sigma(s)$, which exists in a neighborhood of
$\sigma$. Set $k_\sigma=P_4(T)\chi_\sigma$. Integration by parts gives
\[
 \widehat k_\sigma(s)=P_4(s)\widehat\chi_\sigma(s)
                      =1+O((s-\sigma)^5).
\]
Every operation preserves $[a,d]$. We have the direct decomposition
\begin{equation}
 \mathcal B_\sigma
   =\mathbb C k_\sigma\ \oplus\ (T-\sigma)^5\mathcal E.
 \label{eq:ej-Weil-channel}
\end{equation}
To prove both directions, for $g\in\mathcal B_\sigma$ subtract
$\widehat g(\sigma)k_\sigma$. The remainder $g_0$ has all five
Fourier jets of orders zero through four equal to zero at $\sigma$.
The explicit compact inverse is
\[
 f(u)=\frac{i^5e^{i\sigma u}}{4!}
       \int_{-\infty}^u(u-v)^4e^{-i\sigma v}g_0(v)\,dv,
 \qquad (T-\sigma)^5 f=g_0.
\]
Differentiating five times after conjugation by $e^{i\sigma u}$
proves the equation. Above the support interval, expansion of
$(u-v)^4$ gives a polynomial all of whose coefficients are the
five vanishing moments; below it the integral is zero. Smoothness
and all endpoint derivatives follow from the integral. This proves
$f\in\mathcal E$ with no enlarged support. Conversely the Fourier
transform of $(T-\sigma)^5 f$ has all five zero jets, proving
membership. Evaluation at $\sigma$ proves that the sum is direct;
injectivity of $(T-\sigma)^5$ follows because its Fourier product
can vanish identically only if the entire Fourier transform is zero.
This also proves uniqueness of $f$ and of the scalar coefficient.

The exact Weil form on this space is still the original one. If
$g=c k_\sigma+(T-\sigma)^5f$, it is
\begin{align}
 W(g,g)={}&|c|^2W(k_\sigma,k_\sigma)
   +\overline c\,W(k_\sigma,(T-\sigma)^5f)\notag\\
 &+c\,W((T-\sigma)^5f,k_\sigma)
   +W((T-\sigma)^5f,(T-\sigma)^5f).
 \label{eq:ej-Weil-full}
\end{align}
This is sesquilinearity, including both cross terms. It can be
evaluated using the complete pole, gamma and prime formula
\eqref{eq:mag-full-transported-Weil} at $t=0$
\cite{weil1952,alpogeFurman2026}. On a spectral value $s_\rho$,
the last term has the retained factor
$(s_\rho-\sigma)^5\overline{(\overline{s_\rho}-\sigma)^5}$
in the ordering of Section~\ref{sec:retained-weil-algebra};
no assumption that $s_\rho$ is real is used.
The smooth bump generator and the polynomial channel give actual
compact tests, not formal rational functions with unverified inverse domains.

These calculations establish singular diffusion of generic arithmetic
observables at nonzero retained $b$, its exact exceptional jet space,
and the full scalar Weil channels on that space. Neither diffusion
singularity nor the loss of a uniform material bound is itself a
negative value of \eqref{eq:ej-Weil-full}. No such value, or RH
counterexample, is claimed here. The next calculation must use these
actual channels and all cross terms, not replace their pairing by the
positive Euclidean diffusion norm or a signed coefficient.
